% =====================================================================
%  c00.tex ------ 标题页、摘要、译者说明、目录（Lead 撰写，对应原文 p.1--2）
% =====================================================================
\thispagestyle{empty}
\begin{center}
\vspace*{2.2cm}

{\LARGE\bfseries AGT 对应关系慢讲\par}
\vspace{0.6cm}
{\large A Slow Review of the AGT Correspondence\par}
\vspace{0.4cm}
{\normalsize arXiv:2006.14025v3 [hep-th]（2022 年 4 月 6 日）\par}

\vspace{1.4cm}
{\large Bruno Le Floch\par}
\vspace{0.25cm}
{\small CNRS, Laboratoire de Physique Théorique et Hautes Energies,\\
Sorbonne Université, Paris, France\par}

\vspace{1.0cm}
{\small 原文：2022 年 1 月。本文是作者在冬季学校 “YRISW 2020” 所作讲座的扩展写本，
拟收入 J. Phys. A 专刊。\par}

\vspace{1.4cm}
{\large 中译、公式校对与学术批注：本译本\par}
\vspace{0.25cm}
{\small（译本依据 arXiv:2006.14025v3；仅供学习与研究使用）\par}

\vfill
{\small 2026 年 11 月\par}
\end{center}
\clearpage

% ---------------- 摘要（原文 p.1）----------------
\section*{摘要}
\addcontentsline{toc}{section}{摘要}

本讲义从 6d 起源出发，以尽量平缓的方式引入 AGT（Alday--Gaiotto--Tachikawa）对应关系，
并对截至 2020 年初该对应关系众多推广的文献作了广（而不深）的综述。
本文是作者在冬季学校 “YRISW 2020” 所作讲座的扩展写本，拟发表于 J. Phys. A 专刊。

class S 是一大类 4d $\cN=2$ 超对称规范理论（从超 QCD 直至非拉氏量理论），
它们由 6d $\cN=(2,0)$ 超共形场论在黎曼面 $C$ 上作扭曲紧化得到。
这一 6d 构造给出了 class S 理论的 Coulomb 分支与 Seiberg--Witten 几何，将 S 对偶几何化，
并导出了 AGT 对应：class S 理论的许多可观测量等于 2d 共形场论（CFT）的关联函数。
例如，一个 4d $\cN=2$、$SU(2)$ 超共形箭图理论的四球配分函数，等于 Liouville CFT 中主算子的关联函数。

AGT 对应的推广比比皆是：渐近自由规范理论与 Argyres--Douglas 理论对应于非正规（irregular）CFT 算子；
规范群为高阶 rank 的箭图以及 $T_N$ 这样的非拉氏量 tinkertoy 对应于 Toda CFT 关联函数；
非局域算子（Wilson--'t Hooft 圈、面算子、畴壁）则对应于 Verlinde 网络、退化主算子、编织核与融合核，
以及带边界的黎曼面。

\clearpage

% ---------------- 译者说明 ----------------
\section*{译者说明}
\addcontentsline{toc}{section}{译者说明}

\begin{enumerate}[leftmargin=2em,itemsep=4pt]
\item \textbf{底本与页码。}本译本逐页对照 arXiv:2006.14025v3 的 PDF 转写，该 PDF 共 136 页；
      正文为 pp.1--106，pp.106--136 为参考文献（共 1309 条，原文即英文，本译本保留英文原貌）。
      \textbf{原 PDF 页号与原文印刷页码一致（偏移 0）}；源码中以 \verb|% [原文 p.N]| 注释标注页码，
      便于对照原稿核查。
\item \textbf{公式编号。}原文全部 174 条编号公式均与原文一一对应
      （例如式 (3.13) 即原文式 (3.13)）；其中 172 条以 \verb|\tag{}| 手工锁定编号，
      另 2 条为“图形化公式”（式 (1.2)、(1.6)，由原图截取，编号随原图保留）。
      原文中大量公式本就没有编号，译本一律不自行编号。
      正文引用公式写作“式 (3.13)”这样的纯文本，而不使用自动交叉引用（手工 tag 无法自动解析）。
\item \textbf{图与表。}原文有 5 张图（Figure 1--5）、2 张表（Table 1--2）以及一张缩写表。
      原文\textbf{不含任何位图}，全部图形都是矢量绘制；本译本的图与图形化编号公式均按原矢量图形以 400--500 dpi
      渲染截取后重排（因此个别图形公式的字体由原 PDF 直接渲染而来）。
\item \textbf{要点框。}原文每节开头有深灰标题条的要点框（如 “6d viewpoint on 4d $\cN=2$ class S theories”），
      这是原文内容的一部分，本译本以浅灰底框照译。
\item \textbf{批注。}译者在关键处插入 \verb|% [Tutor Note: ...]| 形式的 LaTeX 注释，编译后不进入排版，
      仅供读者在源码中阅读；全部批注汇总见随附文件《译注汇总.md》。
\item \textbf{原文笔误。}原文的印刷错误在译本中\textbf{照原样保留}，另在批注中指出并给出正确形式；
      原文中作者留下的编辑性标记（如 “References welcome”）亦照译保留。
\item \textbf{术语。}以中文学界通行译法为准（见 glossary），首次出现时括注英文；
      人名、模型名与缩写（AGT、class S、CFT、SYM、SQCD 等）保留英文原样。
\item \textbf{强调。}原文的\textbf{粗体}小标题译为中文粗体（\verb|\para{…}|），原文的\emph{斜体强调}
      在译文中统一用粗体表示（中文字体无真斜体，强行倾斜反而降低可读性）。
\item \textbf{版权。}原文版权归作者所有；本译本仅供个人学习与研究使用，不作商业传播。
\end{enumerate}

\vspace{0.6em}
\noindent{\small 说明：原文脚注在译本中由 \LaTeX{} 自动编号，与原文脚注一一对应，但编号数字可能不同。}

\clearpage
\tableofcontents
\clearpage
